\section{Introducción}

El presente reporte desarrolla el análisis de requerimientos y el diseño del
prototipo para el centro de entretenimiento familiar PuellaGame. La afluencia de
clientes durante el periodo vacacional exige comenzar a organizar la
información aun cuando el diseño y la implementación de una base de datos
definitiva requieren un proceso más amplio. Como respuesta inmediata, se plantea
una aplicación que almacena datos persistentes en archivos CSV.

El prototipo administra tres entidades fundamentales: sucursales, premios y
clientes. Para cada una permite registrar, consultar por identificador, editar y
eliminar información mediante un menú de consola. También valida los tipos y
formatos de entrada y maneja condiciones excepcionales para comunicar los
errores sin interrumpir la operación.

Además de describir el contexto, los actores, los requerimientos y la
arquitectura de la solución, el reporte compara el almacenamiento en archivos
con el uso de una base de datos. Esta comparación permite justificar el CSV como
una medida transitoria y determinar la alternativa más conveniente para la
operación futura de PuellaGame.
